% =====================================================================
%  Thème 1 — Conversions & cercle trigonométrique — NIVEAU DIFFICILE
%  Corrigé
% =====================================================================
\documentclass[11pt,a4paper]{article}
\usepackage[utf8]{inputenc}
\usepackage[T1]{fontenc}
\usepackage{lmodern}
\usepackage[french]{babel}
\usepackage{amsmath,amssymb,mathtools}
\usepackage{icomma}
\usepackage{xcolor}
\usepackage{colortbl}
\usepackage{tikz}
\usepackage[most]{tcolorbox}
\usepackage{enumitem}
\usepackage{array}
\usepackage[a4paper,margin=2.2cm]{geometry}
\usetikzlibrary{arrows.meta,calc}

\definecolor{primary}{RGB}{214,51,108}
\definecolor{primarydark}{RGB}{140,20,64}
\definecolor{excol}{RGB}{0,135,90}

\DeclareMathOperator{\tg}{tg}
\DeclareMathOperator{\cotg}{cotg}
\newcommand{\degr}{^{\circ}}
\newcommand{\Z}{\mathbb{Z}}

\newcounter{exo}
\newcommand{\exo}[1]{\medskip\refstepcounter{exo}%
  \noindent{\large\bfseries\color{primarydark}Exercice \theexo}%
  \ \textcolor{primary}{\textbullet}\ \textbf{#1}\par\nopagebreak\smallskip}
\newcommand{\rep}{\textcolor{excol}{$\blacktriangleright$}\ }

\setlength{\parindent}{0pt}
\pagestyle{empty}

\begin{document}

\begin{center}
{\LARGE\bfseries\color{primary}Thème 1 — Conversions \& cercle trigonométrique}\\[2pt]
{\color{primarydark}\textsc{Niveau Difficile — Corrigé détaillé}}
\end{center}
\hrule height 1pt
\medskip

\exo{Le grand angle négatif}
\begin{itemize}[leftmargin=1.4em,itemsep=2pt]
\item[\textbf{(a)}] \rep On ajoute des tours entiers ($2\pi=\frac{12\pi}{6}$) jusqu'à tomber dans
$[0;2\pi[$ :
\[
-\frac{25\pi}{6}+3\times 2\pi=-\frac{25\pi}{6}+\frac{36\pi}{6}=\frac{11\pi}{6}=\beta.
\]
\item[\textbf{(b)}] \rep $\beta=\dfrac{11\pi}{6}\times\dfrac{180}{\pi}=\dfrac{11\times180}{6}=330\degr$ :
c'est le \textbf{quadrant IV}.
\item[\textbf{(c)}] \rep Au quadrant IV : $\cos\beta>0$, $\sin\beta<0$, $\tg\beta<0$.
\item[\textbf{(d)}] \rep Référence $360\degr-330\degr=30\degr$. Comme $\alpha$ et $\beta$ ont les
mêmes nombres trigonométriques :
\[
\cos\alpha=+\cos30\degr=\frac{\sqrt3}{2},\quad
\sin\alpha=-\sin30\degr=-\frac12,\quad
\tg\alpha=-\tg30\degr=-\frac{\sqrt3}{3}.
\]
\end{itemize}

\exo{Mise en situation : la grande roue}
\begin{itemize}[leftmargin=1.4em,itemsep=2pt]
\item[\textbf{(a)}] \rep $210\degr=210\times\dfrac{\pi}{180}=\dfrac{7\pi}{6}$.
\item[\textbf{(b)}] \rep $210\degr$ est entre $180\degr$ et $270\degr$ : \textbf{quadrant III}.
L'ordonnée $y=\sin$ y est \textbf{négative} (la nacelle est en bas).
\item[\textbf{(c)}] \rep Référence $210\degr-180\degr=30\degr$, avec $\cos<0$ et $\sin<0$ :
\[
M=\bigl(\cos210\degr\,;\,\sin210\degr\bigr)=\left(-\frac{\sqrt3}{2}\,;\,-\frac12\right).
\]
\item[\textbf{(d)}] \rep Total $210\degr+270\degr=480\degr$, ramené : $480\degr-360\degr=120\degr$.
$120\degr$ est au quadrant II ($\cos<0,\ \sin>0$), référence $60\degr$ :
\[
M=\bigl(\cos120\degr\,;\,\sin120\degr\bigr)=\left(-\frac12\,;\,\frac{\sqrt3}{2}\right).
\]
\end{itemize}

\exo{Retrouver un angle à partir d'une valeur}
\begin{itemize}[leftmargin=1.4em,itemsep=2pt]
\item[\textbf{(a)}] \rep $\cos x=-\dfrac12$ : valeur négative $\Rightarrow$ quadrants II et III
(référence $60\degr=\frac{\pi}{3}$). Donc
$x=\pi-\dfrac{\pi}{3}=\dfrac{2\pi}{3}$ \ ou \ $x=\pi+\dfrac{\pi}{3}=\dfrac{4\pi}{3}$.
\item[\textbf{(b)}] \rep $\sin x=\dfrac{\sqrt2}{2}$ : valeur positive $\Rightarrow$ quadrants I et II
(référence $45\degr=\frac{\pi}{4}$). Donc
$x=\dfrac{\pi}{4}$ \ ou \ $x=\pi-\dfrac{\pi}{4}=\dfrac{3\pi}{4}$.
\item[\textbf{(c)}] \rep \textbf{Non} : le cosinus est toujours compris entre $-1$ et $1$, or
$\dfrac32>1$. Aucune solution.
\item[\textbf{(d)}] \rep Placement des quatre solutions :
\begin{center}
\begin{tikzpicture}[scale=1.9,line cap=round,>=Stealth]
  \draw[primary,line width=1pt] (0,0) circle (1);
  \draw[->,black!70] (-1.3,0)--(1.35,0) node[right]{\scriptsize$\cos$};
  \draw[->,black!70] (0,-1.3)--(0,1.35) node[above]{\scriptsize$\sin$};
  \foreach \a/\lab/\pos in {120/\frac{2\pi}{3}/{above left}, 240/\frac{4\pi}{3}/{below left},
      45/\frac{\pi}{4}/{above right}, 135/\frac{3\pi}{4}/{above left}}{
     \fill[primarydark] (\a:1) circle (0.028);
  }
  \node[primarydark,font=\scriptsize,above right] at (45:1) {$\frac{\pi}{4}$};
  \node[primarydark,font=\scriptsize,above left] at (135:1) {$\frac{3\pi}{4}$};
  \node[primarydark,font=\scriptsize,left] at (120:1) {$\frac{2\pi}{3}$};
  \node[primarydark,font=\scriptsize,below left] at (240:1) {$\frac{4\pi}{3}$};
  \draw[black!35,dashed] (-1,{-sin(60)})--(1,{-sin(60)});
  \fill (0,0) circle (0.02);
\end{tikzpicture}
\end{center}
\end{itemize}

\end{document}
